\documentclass[11pt]{article}
\usepackage{fontspec}
\usepackage[french,english]{babel}
\usepackage{amsmath}
\usepackage{amssymb}
\usepackage{amsthm}
\usepackage{mathrsfs}
\usepackage{tikz-cd}
\usepackage{enumitem}

\newcommand{\printedpage}[1]{\marginpar{\footnotesize p.~#1}}
\newcommand{\msp}[1]{\marginpar{\footnotesize [#1]}}
% printed notation
\newcommand{\M}{\mathcal{M}}
\newcommand{\Cat}{\mathrm{Cat}}
\newcommand{\Ho}{\mathrm{Ho}}
\newcommand{\Fl}{\mathrm{Fl}}
\newcommand{\Cof}{\mathrm{Cof}}
\newcommand{\Diag}{\mathrm{Diag}}
\newcommand{\id}{\mathrm{id}}
\newcommand{\deq}{\overset{\text{def}}{=}}
% small tilde printed after a symbol (Q-saturation)
\newcommand{\Qs}{\tilde{\ }}
% arrow with one head on the left and a double head on the right (factorisation)
\newcommand{\fact}{\mathrel{\leftarrow\mkern-11mu\twoheadrightarrow}}
% arrow with half heads (harpoons) at both ends, as printed on p. 23
\newcommand{\harp}{\mathrel{\leftharpoondown\joinrel\relbar\joinrel\rightharpoonup}}

\begin{document}

\printedpage{19}
% head: A. Grothendieck : Dérivateurs, Ch. XIX, Édition des Universités de Montpellier II et Paris VII
% foot: Version 11 février 2019

\msp{page 31}
\textbf{Scholie 2.10.} \textit{Let $(W,C,F)$ be a closed Quillen triple in $\M$. For $I$ in \textup{Cat}, for $(W(I),C\langle I\rangle,F(I))$ to be a closed Quillen triple for $\M(I)$, it is necessary and sufficient that the couples $C\langle I\rangle\longleftrightarrow TF(I)$ and $TC\langle I\rangle\longleftrightarrow F(I)$ be factorizing. Suppose that this condition is satisfied for $I$ and for $I'$, and let $u : I\longrightarrow I'$ be an arrow from $I$ to $I'$. Then the assertions of \textup{2.3. 2°} hold, in particular we find a functor $u^{W}_{!} : \Ho_W(I)\longrightarrow\Ho_W(I')$, left adjoint to $u^{*}_{W}$\footnote[18]{Say also that the computation of the fibers of $u^{W}_{!}(X)$, provided for in Der 5 bis, is OK.}. These results dualize at once (by $\M\longmapsto\M^{\circ}$, $I\longmapsto I^{\circ}$) for the triples $(W(I),C(I),F[I])$, and for the construction of $u^{W}_{*}$, right adjoint to $u^{*}_{W}$. In particular, for the preceding triple to be a Quillen triple (necessarily closed), it is necessary and sufficient that the two couples $C(I)\longleftrightarrow TF[I]$ and $TC(I)\longleftrightarrow F[I]$ be factorizing; and when this condition is satisfied for $I$ and for $I'$ (for a given arrow $u : I\longrightarrow I'$ in \textup{Cat}), then $u^{W}_{*}$ exists and is constructed dually to the construction of $u^{W}_{!}$ made explicit in \textup{2.3. 2°}.}

\msp{page 32}
% a vertical bar is printed in the left margin alongside Corollaire 2.11
\textbf{Corollary 2.11.} \textit{The prederivator $\mathbf{D}_{\M,W}$ defined for a \textup{very special} localizer $W$ in $\M$ is a derivator in the full sense of the term, i.e. it satisfies all the axioms Der 1 to Der 6 (including Der 4 bis and Der 5 bis and Der 6 bis) of derivators.}

\textbf{N.B.} To my own surprise, I have not been able to find a case of a localizer $W$ which defines a derivator on all of Cat (with existence of $u^{W}_{!}$, $u^{W}_{*}$), \textit{except} for $W$ a very special localizer, hence fitting into a very special Quillen triple -- or when one can reduce to this in the standard way. (For example in the case where $\M=\Cat$, $W$ a basic localizer satisfying Loc 6 a) and Loc 8, 8 bis)\footnote[19]{axiom of cocartesian squares (direct version), axiom of direct limits, accessibility.} -- the monomorphisms in $W$ flatly refuse to be stable under cobase change!) I have absolutely no idea whether the spectacular result ``proved'' by \textsc{Anderson} is true or false. It looks almost too good to be true \ldots{}Besides, I know of no special Quillen structure which is not very special -- sorry! % [sic] stable -> stables, agreement with « les monomorphismes »

\msp{page 33}

\section*{3\quad Case of categories with cofibrations, of categories with fibrations}

I recall the axioms of a \textit{cofibration structure} $(W,C)$ on a category $\M$.

\begin{itemize}[leftmargin=4em,labelwidth=3.5em,align=left]
\item[Cof$_0$ :] $W$ is a localizer.
\item[Cof$_1$ :] $C$ contains the isomorphisms, is stable under cobase change, under composition.
\item[Cof$_2$ :] $C\cap W$ is stable under cobase change (and also under composition, and contains the isomorphisms, by virtue of Cof$_0$, Cof$_1$).
\item[Cof$_3$ :] Every arrow $f$ in $\M$ factors as $f=pi$, with $i\in C$, $p\in W$.
\end{itemize}

(Axioms of K. S. \textsc{Brown}, slightly generalized by \textsc{Anderson}). There is an optional ``axiom
\printedpage{20}
% head: A. Grothendieck : Dérivateurs, Ch. XIX, Édition des Universités de Montpellier II et Paris VII
% foot: Version 11 février 2019
of continuity'' or \textit{``transfinite axiom''}, which is the following:

\begin{itemize}[leftmargin=4em,labelwidth=3.5em,align=left]
\item[Cof$_4$ :] $\M$ stable under small filtering $\varinjlim$, and $C$ and $TC = C\cap W$ are stable under transfinite composition.
\end{itemize}

Thus, Cof$_1$, Cof$_2$, Cof$_4$ can be summarized by saying that $C$ and $TC$ are transfinitely Q-saturated on the left.

The cofibration structure is called \textit{special} if it satisfies Cof$_4$, if $\M$ is special, and if $C$ admits a \textit{small} subset $C_0$ which is Q-generating on the left, i.e. such that
\[
C \ \subset\ C_0\Qs\,, \qquad \text{i.e.}\quad \tilde{C}=C_0\Qs. \tag*{(3.1)}
\]

\msp{page 34}
Then, by the factorization theorem 1.1, we have
\[
C \ \fact\ TF \ \deq\ C_* \ ,
\]
and the factorization axiom Cof$_3$ implies (cf. 1.4) $TF\subset W^{\natural}$. More generally, we have:

\textbf{Proposition 3.1.} \textit{Suppose that $\M$ is special (stable under small $\varinjlim$ , and accessible elements), that $(W,C)$ satisfies Cof$_0$, Cof$_1$, and that $C$ is stable under transfinite composition (which, together with Cof$_1$, is stated by saying that $C$ is transfinitely Q-saturated on the left), finally that $C$ admits a small subset $C_0$ Q-generating on the left, so that $C\fact TF\deq C_*$. Then for the factorization axiom Cof$_3$ to be satisfied, it suffices that we have}
\[
TF \ \subset\ W \ , \tag*{(3.1.1)}
\]
\textit{and it is necessary to have $TF\subset W^{\natural}$, hence if $W=W^{\natural}$, the axiom Cof$_3$ is equivalent to} (3.1.1).

The cofibration structure is called \textit{semi-closed} if $C$ is Q-closed on the left, i.e. if we have $C=C\Qs$ (which implies that $C$ is stable under transfinite composition), \textit{closed} if in addition $TC=(TC)\Qs$ (which implies the transfinite axiom Cof$_4$).

\msp{page 35}
I do not return here to the theory of K. S. \textsc{Brown} on categories with cofibrations. It is, to my knowledge, the most economical set of axioms known, for ensuring on the one hand a good grip on $\Ho_W$ (through a corresponding notion of homotopy of maps, and a calculus of fractions in $\Pi\M_{\mathrm{c}}$), and on the other hand the existence of the \textit{exact} sequence of cofibrations, hence a theory of $W$-homotopic cofibers. (\textbf{N.B.} When the structure is \textit{proper,} i.e. $C\subset\Cof_W$, then the existence of this structure amounts simply to saying that ``$W$ admits enough $W$-cofibrations'', i.e. that every $f\in\Fl(\M)$ factors as $pi$, with $i\in\Cof_W$, $p\in W$; and we can therefore take $C=\Cof_W$. In the case where $W$ is stable under small filtering $\varinjlim$, $C=\Cof_W$ and $TC=W^{\mathrm{univ}}$ satisfy in addition the transfinite axiom Cof$_4$.) Finally, K. S. \textsc{Brown} shows that every arrow $i : X\longrightarrow Y$ in $C$, \textit{with $X$, $Y$ cofibrant} (i.e. $\emptyset_{\M}\longrightarrow X$, $\emptyset_{\M}\longrightarrow Y$ in $C$) is in $\Cof_W$ (``Brown's Lemma'').

I am interested here in the question of extending the cofibration structure to a category of the form $\M(I)$, and of using this extension to construct
\msp{page 36}
the functors $u^{W}_{!}$. Let us note already that the theory of K. S. \textsc{Brown}, with the supplements of \textsc{Anderson}\footnote[20]{revised and arranged by me.}, gives essentially this:

\printedpage{21}
% head: A. Grothendieck : Dérivateurs, Ch. XIX, Édition des Universités de Montpellier II et Paris VII
% foot: Version 11 février 2019

\textbf{Proposition 3.2.} \textit{Let $W\subset\Fl(\M)$ be a localizer associated with a cofibration structure $(C,W)$ in $\M$.}
\begin{itemize}[leftmargin=3em,labelwidth=2.5em,align=right]
\item[1°)] \textit{On the category \textup{$\Diag_0$} of \textup{finite} ordered sets, this prederivator satisfies the axioms Der 1 (decomposition), Der 2 (localization), Der 3 (final object), Der 4 bis (existence of $u^{W}_{!}$), Der 5 bis (computation of the fibers of $u^{W}_{!}$), Der 6 bis (left exactness).}
\item[2°)] \textit{This is still true on the larger subcategory \textup{$\Diag_1$} in \textup{Cat}, consisting of \textup{noetherian} ordered sets $I$ such that the topological space associated with $I^{\circ}$ is noetherian (every increasing sequence of objects of $I$, or of open sets of $I^{\circ}$, is stationary).}
\item[3°)] \textit{Finally, if the transfinite axiom Cof$_4$ is satisfied, i.e. if $C$ and $TC$ are stable under transfinite composition, then the conclusion of \textup{1°} remains valid on the full subcategory \textup{$\Diag_2$} of \textup{Cat} consisting of ordered sets which are merely noetherian.}
\end{itemize}

\textbf{Corollary 3.3.} \textit{Suppose that we have Cof$_3$ in the stronger form of a factorization
\msp{page 37}
$f=pi$ \textup{functorial} in $f$, which is the case in particular if the cofibration structure is special. Then for every $I$ in \textup{Cat}, $\M(I)$ equipped with $(W(I),C(I))$ is a category with cofibrations, which satisfies the transfinite axiom Cof$_4$ if $(W,C)$ satisfies it. Thus we can apply to it the results of \textup{3.2}, which imply in particular that the axioms Der 1, 2, 3 ( ?) hold on all of \textup{Cat}, and that $u^{W}_{!}$ exists when $u$ is of the form $\id_I\times v$, with $v : J\longrightarrow K$ an arrow in \textup{$\Diag_1$} (resp. in \textup{$\Diag_2$}, if the transfinite axiom is satisfied).}

\textbf{3.4.} The key to the theory is in the extension of the cofibration structure $(W,C)$ from $\M$ to a category $\M(I)$, where $I$ is an ordered set, as $(W(I),C_I)$. The main facts are stated in XVII, in the theorems p. 26-27, p. 50, p. 60. It is when $I$ is \textit{noetherian,} and we assume either that $\mathrm{Esp}(I^{\circ})$ is also noetherian (e.g. $I$ finite), or that $(W,C)$ satisfies the transfinite axiom Cof$_4$, that we prove that $(W(I),C_I)$ is also a cofibration structure on $\M(I)$.

\msp{page 38}
An important fact for the construction of the functors $u^{W}_{!}$, for $u : I\longrightarrow I'$ an arrow between noetherian ordered categories, is the relation
\[
C_I\cap W(I) \ = \ (TC)_I \ , \tag*{(3.4.1)}
\]
from which it follows that we have, for $u$ as above,
\[
u_!(C_I\cap W(I)) \ \subset\ C_{I'}\cap W(I') \ , \tag*{(3.4.2)}
\]
an essential ingredient of the construction of the functor $u^{W}_{!}$. To extend this construction to the case of a $\pm$ arbitrary arrow in Cat, there arises therefore the problem of extending the definition of $C_I$, and the formulas (3.4.1) and (3.4.2), to the case of $\pm$ arbitrary categories in Cat. Now we have proved:

\textbf{Proposition 3.5.} \textit{If $\Phi=\Psi^*$ is a left \textup{Q-closed} subset of \textup{$\Fl(\M)$}, and $I$ a noetherian ordered set, $\Phi_I$ the extension of $\Phi$ to a subset of \textup{$\Fl(\M(I))$} defined by \textup{\textsc{Anderson}}, then we have}
\[
\Phi_I \ = \ \Psi(I)^* \ = \ v_!(\Phi(I_0))\Qs \ , \tag*{(3.5.1)}
\]
\textit{where $v : I_0\longrightarrow I$ is the canonical functor from the discrete category defined by \textup{$\mathrm{Ob}\,I$}.}

\printedpage{22}
% head: A. Grothendieck : Dérivateurs, Ch. XIX, Édition des Universités de Montpellier II et Paris VII
% foot: Version 11 février 2019

This is (with different notations) lemma 6.2.5 (p. 107) of XVII, whose proof I do not repeat here.

This gives the key to a definition of $\Phi_I$ as $\Phi\langle I\rangle$ (with the notations of paragraphs 1, 2),
\msp{page 39}
for arbitrary $I$ in Cat,
\[
\Phi_I \ \deq\ \Phi\langle I\rangle \ \deq\ \Psi(I)^* \ = \ v_!(\Phi(I_0))\Qs \ , \tag*{(3.5.2)}
\]
where $\Psi=\Phi_*$, and where we assume $\Phi$ Q-closed on the left, i.e. $\Phi=\Psi^*$. In the case where we assume only that $\Phi$ is transfinitely Q-saturated, without more, we would like to set
\[
\Phi_I \ \deq\ \text{transfinite Q-saturation of } v_!(\Phi(I_0)) \ , \tag*{(3.5.3)}
\]
and I presume that for $I$ ordered noetherian, this coincides with \textsc{Anderson}'s definition. But I will restrict myself in what follows to the case where $C=\Phi$ is Q-closed on the left.

Let us then return to the case of a cofibration structure $(W,C)$ in $\M$, assuming $C$ and $TC=C\cap W$ Q-closed on the left. Let us define for every $I$ in Cat $C_I=C\langle I\rangle$ as in (3.5.2). We then have, for $u : I\longrightarrow I'$ an arrow in Cat,
\[
u_!(C\langle I\rangle) \ \subset\ C\langle I'\rangle \ , \tag*{(3.5.4)}
\]
but we have little reason here to hope to have (3.4.2) as well, without having the formula
\[
C\langle I\rangle\cap W(I) \ \overset{?}{=}\ (TC)\langle I\rangle \qquad\qquad \text{\footnote[21]{\textbf{N.B.} We have the inclusion $\supset$, since $v^*v_!TC(I_0)\subset TC(I_0)$.}}. \tag*{(3.5.5)}
\]
We have proved that this formula is valid if $I$ is ordered noetherian (XVII th. 2.1, p. 26-27). This
\msp{page 40}
brings us to the

\textbf{Question 3.6.} \textit{Let $\M$ be a category stable under small limits, and let $(W,C)$ be a cofibration structure in $\M$ such that $C$ and $TC=C\cap W$ are Q-closed on the left. Let $I$ be an object of \textup{Cat}, do we then have}
\[
C\langle I\rangle\cap W(I) \ = \ (TC)\langle I\rangle \text{ ?} \tag*{(3.6.1)}
\]

We saw in the preceding paragraph that this is so if $(W,C)$ fits into a closed Quillen triple $(W,C,F=(TC)_*)$, and when we assume in addition that the couples
\[
C\langle I\rangle \ \longleftrightarrow\ TF(I) \ , \qquad TC\langle I\rangle \ \longleftrightarrow\ F(I)
\]
are factorizing, in which case $(W(I),C\langle I\rangle,F(I))$ is a closed Quillen structure in $\M(I)$. Besides, $(W,C,F)$ is a Quillen structure (necessarily closed) if and only if the couples $C\longleftrightarrow \underbrace{TF}_{\deq C_*}$, $TC\longleftrightarrow \underbrace{\ F\ }_{\deq (TC)_*}$ themselves are factorizing, and if % [sic] factorisant -> factorisants, agreement with « les couples »
\printedpage{23}
% head: A. Grothendieck : Dérivateurs, Ch. XIX, Édition des Universités de Montpellier II et Paris VII
% foot: Version 11 février 2019
in addition $TF\subset W$ (which is automatic if $W=W^{\natural}$, for we will have in any case, by Cof$_4$, $TF\subset W^{\natural}$).

But can we get by with less, to be able to conclude (3.6.1)? Would it not suffice to know that $TC\langle I\rangle\longleftrightarrow F(I)$ is factorizing (which has little chance of being true for an $I$ $\pm$ general, if we do not already have $TC\fact F$, factorization in $\M$ itself). This is indeed the case:

\msp{page 41}
\textbf{Proposition 3.7.} \textit{Under the general conditions of \textup{(3.6)} on $\M$, $W$, $C$, $I$, suppose that the couple $TC\langle I\rangle\longleftrightarrow F(I)$ (where $F\deq (TC)_*$) is factorizing, i.e. $TC\langle I\rangle\fact F(I)$. Then the formula \textup{(3.6.1)} is true.}

Indeed, if $f\in C\langle I\rangle\cap W(I)$, we factor $f$ as $pi$, with $i\in TC\langle I\rangle$, $p\in F(I)$. Since $f,i\in W(I)$, we will have $p\in W(I)$, i.e. $p\in TF(I)$, and because of $\underbrace{C\langle I\rangle}_{\ni\, f}\harp\underbrace{TF(I)}_{\ni\, p}$, there exists $p : Y\longrightarrow X'$ such that $pq=\id_Y$, $qf=i$, % reading resolved: DER_XIX_corrections.csv; % [sic] p : Y -> q : Y, the arrow Y -> X' is q (see the diagram and pq = id_Y)
\[
\begin{tikzcd}[column sep=small, row sep=large]
X \arrow[rr, "f\in C\langle I\rangle\cap W(I)"] \arrow[dr, "i\in TC\langle I\rangle"'] & & Y \arrow[dl, dotted, shift right=1ex, "q"'] \\
& X' \arrow[ur, "p\in F(I)"'] \ , &
\end{tikzcd}
\]
so $f$ is a direct factor of $i\in TC\langle I\rangle$. Since $TC\langle I\rangle$ is stable under direct factors, we win.

\textbf{Corollary 3.8.} \textit{Let $(W,C,F)$ be a closed Quillen structure in $\M$, and suppose that $TC=C\cap W$ is a Q-special subset of $\M$, i.e. the couple $TC\fact F$ a special couple (i.e. $\M$ is \textup{special}, and there exists a small subset $TC_0$ Q-generating on the left for $TC$). Then for every $I$ in \textup{Cat}, we have}
\[
C\langle I\rangle\cap W(I) \ = \ TC\langle I\rangle \ . \tag*{(3.8.1)}
\]

\msp{page 42}
Apparently, we have not had to assume that $C$ too is a Q-special subset of $\Fl(\M)$. But as we need the factorization for $C\longleftrightarrow TF$ (to have at least a cofibration structure on $\M$), we have hardly any way of establishing it without assuming $C$ special. Not to mention that to be able to make (3.8.1) explicit for the construction of $u^{W}_{!}$, we need to know that $(W(I),C\langle I\rangle)$ does indeed define a cofibration structure on $\M(I)$, which can hardly be done (it seems to me) without having the factorization for $C\langle I\rangle\longleftrightarrow TF(I)$ as well; something I do not manage to establish for a general $I$ in Cat, except by means of the hypothesis of specialness on $C$ as well.

\textbf{3.9.} To get to the bottom of it, we would need a few concrete and illuminating examples, of the following type:

a) \texttt{[There is no b).]} Closed Quillen structure $(W,C,F)$ on $\M$, such that for the associated prederivator, $u^{W}_{!}$ does not always exist, or $u^{W}_{*}$ does not exist by always. I am thinking in particular of the case of an arrow $u : I\longrightarrow I'$ with $I'=e$ (presuming that if in this case, $u^{W}_{!}$ resp. $u^{W}_{*}$ exists, then it exists in all cases). Such an example would be even more % [sic] n'existe par toujours -> n'existe pas toujours, typing error

\end{document}
